\section{Экспериментальная проверка}

\subsection{Стенд, нагрузки и подтверждение корректности}

Экспериментальное исследование проведено на Linux-машине GitHub Actions с CPU, Python~3.11, PyTorch~2.5.1 и Ray~2.40.0. Помимо ранее защищённого прототипа НИР, создан \emph{отдельный воспроизводимый экспериментальный модуль}, который реализует монолитный базовый режим B0 и распределённые варианты D0--D3. На одном физическом вычислительном экземпляре Ray создано три независимых процесса-исполнителя. Это не эквивалентно трём физическим серверам и не заменяет сетевые эксперименты исходного НИР.

Сформирован набор из 144 конфигураций: четыре категории исходных текстов, по три авторских текстовых шаблона для каждой категории, четыре длины входной последовательности (16, 64, 128, 256 токенов) и три длины генерации (8, 16, 32 токена). Таким образом, набор содержит \emph{12 авторских исходных текстов и 144 сочетания параметров}, а не 144 независимых естественных запроса. Для каждого сценария фактические входные токены формируются токенизатором GPT-2 и сохраняются вместе с параметрами генерации.

На \code{sshleifer/tiny-gpt2} выполнены $144\cdot 5\cdot 3=2160$ измерений после отдельного прогрева: по три повтора каждого режима на каждом сценарии. Для каждого распределённого запуска проверялось точное совпадение последовательности генерируемых идентификаторов с монолитным эталоном. Результат: \textbf{2160 совпадений из 2160}; ошибки и несовпадения отсутствуют в журнале. В данных сохранены длительности, тип нагрузки, режим, количество возвращённых байтов и параметры исполнения. Успешная сборка отчёта и тестов доступна в открытом репозитории исследовательского модуля.

\subsection{Варианты обработки выходных логитов}

\begin{table*}[t]
\centering\small
\caption{Длительность инференса на tiny-gpt2: полная серия, 144 сценария и три повтора}
\label{tab:full-results}
\begin{tabularx}{\textwidth}{@{}p{1.3cm}p{4.1cm}rrX@{}}
\toprule
\textbf{Режим} & \textbf{Обработка выхода} & \textbf{Медиана, мс} & \textbf{90-й перцентиль, мс} & \textbf{Медиана сырого выходного массива}\\
\midrule
B0 & Монолит, полная выходная проекция & 48,65 & 382,49 & Не передаётся через Ray\\
D0 & Полные логиты всех позиций & 178,48 & 530,55 & 323\,655\,080 байт\\
D1 & Возврат логитов последней позиции & 134,06 & 470,13 & 3\,216\,448 байт\\
D2 & Вычисление логитов только последней позиции & 81,95 & 169,05 & 3\,216\,448 байт\\
D3 & Локальный \code{argmax}, возврат ID & 71,36 & 148,57 & 128 байт\\
\bottomrule
\end{tabularx}
\end{table*}

Результаты представлены в табл.~\ref{tab:full-results}. По медиане парных отношений задержек для одинаковых сценариев и повторов отношение D0/D3 составило около $1{,}86$. Эффект включает \emph{два} изменения: уменьшение передаваемого сообщения и устранение вычисления ненужных логитов. Поэтому выводить из него сетевое ускорение было бы некорректно. В качестве полезного объёма фиксировалось число байтов возвращённых массивов (\code{nbytes}); фактический сетевой трафик, стоимость сериализации и издержки RPC отдельно не измерялись.

\subsection{Контроль с равноценно оптимизированным монолитом}

Для устранения методологической несопоставимости B0 и D3 введён дополнительный \emph{контрольный режим B1}: монолитная модель так же вычисляет выходную линейную проекцию только для последней позиции. В отдельной серии выполнены $144\cdot 3\cdot 3=1296$ запусков B0, B1 и D3; все результаты генерации совпали с эталоном. Медианные задержки в \emph{этой одной серии} составили: B0~--- 60,48 мс, B1~--- 6,15 мс, D3~--- 43,95 мс. Медианное парное отношение D3/B1 составило $7{,}24$.

Таким образом, хотя D3 быстрее исходного распределённого D0, на крайне малой модели и одном хосте он существенно \emph{медленнее} равноценно оптимизированного монолита. Это показывает значимость фиксированных расходов Ray относительно малых вычислительных затрат и исключает некорректную интерпретацию сравнения B0--D3 как доказательства пользы распределения самой модели.

\subsection{Шестислойная модель и размещение по динамическому программированию}

Дополнительно проверена \code{distilbert/distilgpt2}, содержащая шесть Transformer-блоков. При равномерном размещении блоки распределены по исполнителям как $[0,2)$, $[2,4)$, $[4,6)$; во всех 20 проверках пяти режимов на четырёх сценариях получены те же токены, что у монолита. Вычисления действительно исполнялись тремя разными процессами.

Для исследования алгоритма выбора границ построен \emph{синтетический} профиль неоднородных скоростей, памяти и стоимости границ. Динамическое программирование получило разбиение $(0,1,2,6)$ против равномерного $(0,2,4,6)$. Оба распределения затем реально исполнены с шестислойной DistilGPT2, по два повторения каждого из пяти режимов на четырёх сценариях: всего 40 корректных сравнений на каждое распределение. Математическая оптимальность подтверждена сравнением с полным перебором для малых конфигураций.

Время режима D3 в этих пилотных сериях имело медиану около 908,37 мс для выбранного алгоритмом разбиения и 919,96 мс для равномерного. Разница мала и не может служить доказательством выигрыша: \emph{профиль для DP не был измерен на неоднородных физических узлах}, а обе серии исполнялись на общей вычислительной машине. Тем самым подтверждена интеграция математического выбора границ в работающий инференс, но не оптимальность на реальном кластере.

\subsection{Ограничения экспериментальной интерпретации}

Полученные числа относятся к конкретному Linux CPU-стенду GitHub Actions; их нельзя переносить на MacBook с Apple M1 Pro, GPU-системы или физически независимые серверы. Набор из 144 конфигураций основан на 12 текстовых шаблонах и не обеспечивает статистической репрезентативности произвольных пользовательских запросов. Эксперименты выполнялись без KV-кэша, пакетной обработки и реального многоузлового сетевого профилирования. Для практического подтверждения оптимального разбиения на неоднородном оборудовании необходимы инструментально измеренные профили вычислений и передачи данных на каждом узле, которые в данном исследовании пока не получены.

